% !PRINTED 682
\weberpage{682}
Es ist aber nach dem Multiplicationssatze der Determinanten

\weq{3}{\begin{vmatrix} \lambda_{1,1}, & \cdots, & \lambda_{\nu-1,1} \\ \cdot & \cdot & \cdot \\ \lambda_{1,\nu-1}, & \cdots, & \lambda_{\nu-1,\nu-1} \end{vmatrix}
=
\begin{vmatrix} \xi_{1,1}, & \cdots, & \xi_{1,\nu-1} \\ \cdot & \cdot & \cdot \\ \xi_{\nu-1,1}, & \cdots, & \xi_{\nu-1,\nu-1} \end{vmatrix}
\begin{vmatrix} l_{1,1}, & \cdots, & l_{\nu-1,1} \\ \cdot & \cdot & \cdot \\ l_{1,\nu-1}, & \cdots, & l_{\nu-1,\nu-1} \end{vmatrix},}

oder, wenn wir

\[
\Lambda_{\nu-1} = \Sigma \pm \lambda_{1,1}\lambda_{2,2} \cdots \lambda_{\nu-1,\nu-1}
\]

setzen, und beachten, dass die Determinante

\[
\Sigma \pm \xi_{1,1}\xi_{2,2} \cdots \xi_{\nu-1,\nu-1}
\]

eine rationale Zahl mit dem Nenner $m^{\nu-1}$ ist:

\weq{4}{\Lambda_{\nu-1} = \frac{a}{m^{\nu-1}} \, L(\varepsilon_1, \varepsilon_2, \ldots, \varepsilon_{\nu-1}),}

worin $a$ eine ganze rationale Zahl ist.

Daraus ergiebt sich, dass die neu eingeführten Einheiten
immer und nur dann ein unabhängiges System bilden, wenn die
Determinante der $\xi_{i,k}$, d. h. die Zahl $a$, von Null verschieden ist.

Nun giebt es unter einer endlichen oder unendlichen Anzahl
nicht verschwindender rationaler Brüche mit demselben Nenner
immer einen dem absoluten Werthe nach kleinsten, und folglich
können wir nach (4) das neue System unabhängiger Einheiten
so wählen, dass sein Regulator $\pm \Lambda_{\nu-1}$ \gesperrt{so klein als möglich}
wird. Ein solches System nennen wir ein \gesperrt{Fundamentalsystem
von Einheiten}, und den Minimalwerth des Regulators selbst,
also den Regulator eines Fundamentalsystemes, nennen wir den
\gesperrt{Regulator des Körpers}. Wir nehmen jetzt an, dass unser
System $\varepsilon_1, \varepsilon_2, \ldots, \varepsilon_{\nu-1}$ selbst ein Fundamentalsystem sei.

Unter dieser Voraussetzung lässt sich beweisen, dass
alle Exponenten von Einheiten \gesperrt{ganze Zahlen} sein müssen.

Wenn wir nämlich annehmen, es existire eine Einheit $\varepsilon$,
deren nach den Formeln \S.\,186, (5) bestimmte Exponenten nicht
alle ganze Zahlen sind, so giebt es nach \S.\,186, 4) auch eine
Einheit $\varepsilon_0$, deren Exponenten echte Brüche sind, die nicht alle
verschwinden.

Verstehen wir unter $\lambda_1, \lambda_2, \ldots, \lambda_{\nu-1}$ in den Formeln (5),
\S.\,186 das System der conjugirten Logarithmen von $\varepsilon_0$, so ergiebt
 sich

\[
L(\varepsilon_0, \varepsilon_2, \ldots, \varepsilon_{\nu-1}) = \xi_1\, L(\varepsilon_1, \varepsilon_2, \ldots, \varepsilon_{\nu-1}),
\]
%CONT>>
